
\documentclass{article}
\usepackage{colortbl}
\usepackage{makecell}
\usepackage{multirow}
\usepackage{supertabular}

\begin{document}

\newcounter{utterance}

\centering \large Interaction Transcript for game `eqbench', experiment `eq\_bench', episode 157 with reflexion{-}llama{-}3.1{-}8b{-}instruct.
\vspace{24pt}

{ \footnotesize  \setcounter{utterance}{1}
\setlength{\tabcolsep}{0pt}
\begin{supertabular}{c@{$\;$}|p{.15\linewidth}@{}p{.15\linewidth}p{.15\linewidth}p{.15\linewidth}p{.15\linewidth}p{.15\linewidth}}
    \# & \multicolumn{2}{c}{Player} && \multicolumn{2}{c}{Game Master} \\
    \hline

    \theutterance \stepcounter{utterance}  
    & & & \multicolumn{4}{p{0.6\linewidth}}{
        \cellcolor[rgb]{0.9,0.9,0.9}{
            \makecell[{{p{\linewidth}}}]{
                \texttt{\tiny{[P1$\langle$GM]}}
                \texttt{You have attempted this same task before and did not succeed. These are the plans you wrote after those attempts.} \\
\texttt{Trial \#1: New plan: I overestimated the likelihood of Eilis feeling flattered by Remy's statement, given their previous exchange. I also underestimated the intensity of Eilis's skepticism, as she has been consistently critical of Remy's motivations. To improve my assessment, I will pay closer attention to the tone and language used in the dialogue, particularly in the final exchange. I will also consider the context of the conversation and the characters' past interactions. Next time, I will give more weight to Eilis's skeptical response and consider the possibility that she may be feeling a mix of emotions, rather than a single dominant one.} \\
\\ 
\texttt{Trial \#2: New plan: I overestimated the likelihood of Eilis feeling flattered by Remy's statement, given their previous exchange. I also underestimated the intensity of Eilis's skepticism, as she has been consistently critical of Remy's motivations. To improve my assessment, I will pay closer attention to the tone and language used in the dialogue, particularly in the final exchange. I will also consider the context of the conversation and the characters' past interactions. Next time, I will give more weight to Eilis's skeptical response and consider the possibility that she may be feeling a mix of emotions, rather than a single dominant one.} \\
\\ 
\texttt{To achieve this, I will:} \\
\\ 
\texttt{1. Analyze the dialogue for subtle cues, such as tone, language, and body language, to better understand the emotional dynamics between Eilis and Remy.} \\
\texttt{2. Consider the context of the conversation, including the characters' past interactions and the stakes involved, to inform my emotional assessment.} \\
\texttt{3. Be more nuanced in my emotional scoring, allowing for the possibility of multiple emotions and varying intensities.} \\
\texttt{4. Pay closer attention to the final exchange, where Remy's statement and Eilis's response may provide valuable insight into her emotional state.} \\
\\ 
\texttt{By implementing these changes, I aim to improve my emotional assessment and provide a more accurate prediction of Eilis's emotional response.} \\
\\ 
\\ 
\\ 
\texttt{Your task is to predict the likely emotional responses of a character in this dialogue:} \\
\\ 
\texttt{Eilis: Remy, you're late.} \\
\texttt{Remy: And you're as punctual as ever, Eilis. Still chasing shadows, are we?} \\
\texttt{Eilis: Better shadows than illusions of grandeur, Remy.} \\
\texttt{Remy: Ah, the sting of your wit. I've missed that. } \\
\texttt{Eilis: And I've missed your incessant need to prove something.} \\
\texttt{Remy: Is that so? Well, I suppose we're both here for the same reason then.} \\
\texttt{Eilis: Don't flatter yourself, Remy. Our reasons couldn't be more different.} \\
\texttt{Remy: Oh, do enlighten me, Eilis. } \\
\texttt{Eilis: You're here to feed your ego. I'm here to reclaim what I lost.} \\
\texttt{Remy: And what if I told you I'm here for you, Eilis?} \\
\texttt{{[}End dialogue{]}} \\
\\ 
\texttt{At the end of this dialogue, Eilis would feel...} \\
\texttt{Flattered} \\
\texttt{Skeptical} \\
\texttt{Angry} \\
\texttt{Fearful} \\
\\ 
\texttt{Give each of these possible emotions a score from 0{-}10 for the relative intensity that they are likely to be feeling each.} \\
\\ 
\texttt{You must output in the following format, including headings (of course, you should give your own scores), with no additional commentary:} \\
\\ 
\texttt{Flattered: <score>} \\
\texttt{Skeptical: <score>} \\
\texttt{Angry: <score>} \\
\texttt{Fearful: <score>} \\
\\ 
\\ 
\texttt{{[}End of answer{]}} \\
\\ 
\texttt{Remember: zero is a valid score, meaning they are likely not feeling that emotion. You must score at least one emotion > 0.} \\
\\ 
\texttt{Your answer:} \\
            }
        }
    }
     \\ \\

    \theutterance \stepcounter{utterance}  
    & \multicolumn{4}{p{0.6\linewidth}}{
        \cellcolor[rgb]{0.561,0.737,0.561}{
            \makecell[{{p{\linewidth}}}]{
                \texttt{\tiny{[P1$\rangle$GM]}}
                \texttt{Flattered: 0} \\
\texttt{Skeptical: 9} \\
\texttt{Angry: 6} \\
\texttt{Fearful: 0} \\
            }
        }
    }
    & & \\ \\

    \theutterance \stepcounter{utterance}  
    & & & \multicolumn{2}{p{0.3\linewidth}}{
        \cellcolor[rgb]{0.9,0.9,0.9}{
            \makecell[{{p{\linewidth}}}]{
                \texttt{\tiny{[GM$|$GM]}}
                \texttt{Flattered: 0} \\
\texttt{Skeptical: 9} \\
\texttt{Angry: 6} \\
\texttt{Fearful: 0} \\
            }
        }
    }
    & & \\ \\

    \theutterance \stepcounter{utterance}  
    & & & \multicolumn{2}{p{0.3\linewidth}}{
        \cellcolor[rgb]{0.9,0.9,0.9}{
            \makecell[{{p{\linewidth}}}]{
                \texttt{\tiny{[GM$|$GM]}}
                \texttt{Flattered: 2} \\
\texttt{Skeptical: 8} \\
\texttt{Angry: 4} \\
\texttt{Fearful: 0} \\
            }
        }
    }
    & & \\ \\

    \theutterance \stepcounter{utterance}  
    & & & \multicolumn{2}{p{0.3\linewidth}}{
        \cellcolor[rgb]{0.9,0.9,0.9}{
            \makecell[{{p{\linewidth}}}]{
                \texttt{\tiny{[GM$|$GM]}}
                \texttt{game\_result = LOSE} \\
            }
        }
    }
    & & \\ \\

    \theutterance \stepcounter{utterance}  
    & & \multicolumn{4}{p{0.6\linewidth}}{
        \cellcolor[rgb]{0.95,0.95,0.95}{
            \makecell[{{p{\linewidth}}}]{
                \texttt{\tiny{[P1$\leftrightarrow$P]}}
                \texttt{Post{-}episode Reflexion (not sent to GM):} \\
\texttt{New plan: I noticed that I've been consistently overestimating the likelihood of Eilis feeling flattered by Remy's statement, despite their previous critical exchange. To improve my assessment, I will pay closer attention to the tone and language used in the dialogue, particularly in the final exchange. I will also consider the context of the conversation and the characters' past interactions. Next time, I will give more weight to Eilis's skeptical response and consider the possibility that she may be feeling a mix of emotions, rather than a single dominant one.} \\
\\ 
\texttt{To achieve this, I will:} \\
\\ 
\texttt{1. Analyze the dialogue for subtle cues, such as tone, language, and body language, to better understand the emotional dynamics between Eilis and Remy.} \\
\texttt{2. Consider the context of the conversation, including the characters' past interactions and the stakes involved, to inform my emotional assessment.} \\
\texttt{3. Be more nuanced in my emotional scoring, allowing for the possibility of multiple emotions and varying intensities.} \\
\texttt{4. Pay closer attention to the final exchange, where Remy's statement and Eilis's response may provide valuable insight into her emotional state.} \\
\\ 
\texttt{However, I also noticed that I've been repeating the same plan in both trials, which suggests that I may not be fully incorporating the lessons learned from the previous attempts. To break this cycle, I will:} \\
\\ 
\texttt{5. Take a step back and re{-}evaluate the dialogue from a fresh perspective, looking for new insights and perspectives that may have been overlooked in previous attempts.} \\
\texttt{6. Consider alternative emotional responses that may be more plausible given the context and tone of the dialogue.} \\
\texttt{7. Develop a more detailed and specific plan for how I will incorporate these new insights and perspectives into my emotional assessment.} \\
\\ 
\texttt{By taking a more nuanced and iterative approach, I aim to improve my emotional assessment and provide a more accurate prediction of Eilis's emotional response.} \\
            }
        }
    }
    & \\ \\

\end{supertabular}
}

\end{document}
